% The v--t diagram for constant acceleration: the line whose height is
% v_f = v_0 + a(Delta t), and the area under it, which is the displacement.
%
% The elapsed interval is lettered Delta t rather than t, so that it is not
% read as the coordinate the axis is already called t. The final speed follows
% it to v_f: a v_t here would put that same bare t back in the picture, in the
% one place a reader is looking for the interval.
%
% Pulled in with \tikzfig{cha-vt-diagram-main.tex}. A fragment, not a document
% -- one tikzpicture, no preamble -- so it inherits the fonts of whatever
% \inputs it. The library loads are idempotent and keep the fragment
% self-contained.
%
% Why the trapezoid is split the way it is: the whole point of the diagram is
% that the area under the line reads as a rectangle v_0 x (Delta t) stacked
% under a triangle of height a(Delta t), which is where the two terms of
% Delta x = v_0(Delta t) + (1/2)a(Delta t)^2 come from. So the v_0 level line
% runs the full width -- it is the lid of the rectangle, not merely a tick --
% and the dimensions at x = \xdim are stacked v_0 then a(Delta t), in the
% order the terms appear in the result.
%
% The level lines and the t axis are continued to the right as dashes so the
% overall v_f dimension has something to measure between. That continuation is
% a construction line rather than part of the graph, which is what the change
% of dash pattern says.
%
% Everything in the picture is derived from the two groups of constants below:
% three that say what the motion is, and four that say how loosely it is set.
% Nothing else is a bare number, so the figure can be re-proportioned or
% re-lettered without hunting for offsets that happened to look right.

% Kept outside the picture because the style below expands them: the barb's
% included angle is what a maintainer thinks in, but Straight Barb has no key
% for it, so the width is computed from it rather than typed as a decimal.
\pgfmathsetmacro{\barbAng}{60}                              % degrees
\pgfmathsetmacro{\barbLen}{6}                               % pt
\pgfmathsetmacro{\barbWid}{2*\barbLen*tan(\barbAng/2)}      % pt

% Needs the arrows.meta, calc and patterns.meta libraries. They are NOT loaded
% here: \usetikzlibrary inside a figure environment sets its "loaded" flag
% globally but its definitions only for that group, which silently disarms the
% next figure's load. gwbook.cls and _wrapper.tex load them in the preamble
% instead -- add to those two if this figure ever needs another.
\begin{tikzpicture}[
    % The figsize knob: distances scale, lettering and hatching do not.
    scale=1.15,
    >={Latex},
    axis/.style  = {->, line width=0.9pt},
    % The graph itself is the heaviest stroke in the picture; everything else
    % is there to measure it.
    graph/.style = {line width=1.3pt},
    % The ordinate at Delta t closes the shaded area but asserts nothing about
    % the motion -- it is where we chose to stop, not something the diagram is
    % claiming -- so it is drawn at a construction weight. At the graph's 1.3pt
    % it read as a second heavy line and made the trapezoid look like a box.
    guide/.style = {line width=0.7pt},
    % Construction: levels carried across to the dimensions on the right.
    % Long dashes, no dots: at this line width a dot reads as a blemish rather
    % than as part of the pattern. The gap is held at 40% of the dash, which is
    % what keeps the line scanning as one level rather than a row of ticks --
    % so tighten by shortening both together, never by closing the gap alone.
    % Measured at this 0.6pt width: 2.5pt/1pt is tighter still, 2pt/0.8pt is
    % about the floor, and by 1.5pt/0.7pt the dashes silt up and read solid.
    level/.style = {dash pattern=on 3pt off 1.2pt, line width=0.6pt},
    dim/.style   = {<->, line width=0.7pt},
    % The hatch reads as "area" in black and white, which a grey fill does not.
    % Lines[] is the stock north east lines hatch with its spacing exposed.
    % Stock sits near 2pt, which fills in to a grey wash at print size; 3.5pt
    % keeps the texture reading as separate strokes without thinning out into
    % a few stray diagonals. distance is the knob -- raise it for sparser.
    area/.style  = {pattern={Lines[angle=45, distance=3.5pt, line width=0.45pt]},
                    pattern color=black!60},
    % The pointer to the result is the one informal mark in the figure -- a
    % pen gesture rather than a measurement -- so it gets a round cap and a
    % single unbroken sweep instead of the dimension arrows' square ends.
    % Straight Barb is an open V of two strokes, which is the head a pen
    % actually makes; the filled triangles (Latex, Stealth) are draughtsman's
    % marks and belong on the dimensions, where they are already doing the job
    % of pointing at an exact height.
    soft/.style  = {->, >={Straight Barb[round, length=\barbLen pt,
                                         width=\barbWid pt]},
                    line width=0.7pt, line cap=round, line join=round},
  ]

  % --- what the motion is ----------------------------------------------------
  % Three chosen numbers; v_f is derived, so the picture cannot end up
  % asserting a height that disagrees with its own slope.
  \pgfmathsetmacro{\vz}{2.0}        % v_0, the intercept
  \pgfmathsetmacro{\T}{5.0}         % Delta t, the interval being marked off
  \pgfmathsetmacro{\rise}{1.3}      % a(Delta t), the speed gained over it
  \pgfmathsetmacro{\vf}{\vz+\rise}  % v_f, derived, never typed

  % --- how loosely it is set -------------------------------------------------
  % \gap   the one clearance unit: between any two columns of the right-hand
  %        margin, and between the v axis and the dimension left of it
  % \lblw  room for the Delta v label, which is what holds the overall v_f
  %        dimension clear of it. The one measured quantity here, and so the
  %        one thing to re-check if that label is ever re-lettered longer.
  %        Measure, do not guess: \sbox the label in gwbook.cls and read \the\wd,
  %        then divide by 1cm and by the picture's scale to get these units. It
  %        must fit the BOOK's setting, which is the wider of the two -- the
  %        standalone export sets the same label about 17% narrower.
  % \stub  how far the t axis and the v_0 level poke out past the left
  %        dimension, so that dimension has something to sit against
  % \runout how far an axis runs past the content it carries, before its tip
  %        (named around \over, which is a TeX primitive and must not be
  %        shadowed in a group where math is typeset)
  \pgfmathsetmacro{\gap}{0.55}
  \pgfmathsetmacro{\lblw}{1.95}   % $\Delta v = a\,\Delta t$: 55.8pt + node padding
  \pgfmathsetmacro{\stub}{0.35}
  \pgfmathsetmacro{\runout}{0.95}

  % The right margin, left to right: the stacked v_0 / a(Delta t) dimension,
  % the room its labels need, then the overall v_f dimension, then the levels
  % run one gap past everything. The t axis stops one gap short of the overall
  % dimension, so its arrow does not read as pointing at it.
  \pgfmathsetmacro{\xleft}{-\gap}
  \pgfmathsetmacro{\xdim}{\T+\gap}
  \pgfmathsetmacro{\xtot}{\xdim+\lblw+\gap}
  \pgfmathsetmacro{\xtip}{\xtot-\gap}
  \pgfmathsetmacro{\xlev}{\xtot+\gap}

  \coordinate (O)  at (0,0);
  \coordinate (A)  at (0,\vz);      % where the motion starts
  \coordinate (B)  at (\T,\vf);     % ... and where it is at the end of Delta t
  \coordinate (Bt) at (\T,0);

  % --- the area --------------------------------------------------------------
  % Filled before anything is drawn, so the axes and the graph sit on top of
  % the hatching rather than under it. Its four sides are the v axis, the t
  % axis, the ordinate at Delta t and the graph -- all drawn below in their own
  % right (the ordinate as a guide, the rest at their own weights), so this
  % path only fills and never strokes.
  \path[area] (O) -- (Bt) -- (B) -- (A) -- cycle;

  % --- axes ------------------------------------------------------------------
  \draw[axis] (\xleft-\stub,0) -- (\xtip,0) node[below right=1pt] {$t$};
  \draw[axis] (0,-\runout) -- (0,\vf+\runout) node[above left=1pt] {$v$};

  % --- the graph -------------------------------------------------------------
  % Extended 15% past B rather than to a typed endpoint: the continuation is
  % then guaranteed collinear with the segment being measured, which is the
  % one thing a reader checks by eye.
  \draw[graph] (A) -- ($(A)!1.15!(B)$);
  \draw[guide] (Bt) -- (B);

  % The slope is the one thing the line means, so it is lettered plainly across
  % the top and tied to the line by a pen stroke, rather than set along the line
  % itself: a tilted label costs the reader more than following an arrow does.
  % Hung from the top of the v axis, so it rises and falls with the picture, and
  % pinned by its left edge one clearance in from the axis. Centring it instead
  % makes the clearance depend on the label's own width -- which is not fixed:
  % the same fragment sets wider in the book's body size than in the standalone
  % export, and a centred label duly reached back over the v axis there.
  \node (slope) [anchor=north west] at (\gap,\vf+\runout)
        {$\dfrac{\Delta v}{\Delta t} = a$ \ (slope $=$ acceleration)};

  % Both ends are pinned to things that do not move: an offset from the label's
  % bottom-left corner, which is fixed because the label is anchored there, and
  % a point ON the segment, named the way the graph itself is. Hanging it off
  % .south instead would tie the stroke to the label's centre, and so to its
  % width -- the one thing that is not fixed between the book and the export.
  %
  % It leaves from the label's own baseline -- .base west is the text baseline
  % at the pinned left edge, so the stroke starts under the lettering the way a
  % pen would, not from the corner of an invisible box -- and runs down-right at
  % about 45 degrees onto the middle of the segment. Leaving steeper than it
  % arrives is what bows it; a vertical drop with a hook at the end reads as a
  % different gesture entirely.
  % shorten keeps the barb off the heavy stroke rather than merging into it.
  % Placed a fraction of the way along the label's own baseline, so it sits
  % under the "=" at any body size: the label scales as a whole, so a
  % proportion of its width tracks a given character in it. (That is not the
  % width-dependence the label's own anchoring avoids -- clearance from a fixed
  % object must not follow the width, but a point meant to track the lettering
  % must.)
  %
  % The stroke drops, turns back up, then drops again into the line. That is one
  % cubic Bezier -- all `to[out/in]' ever draws -- and the three phases come from
  % two independent things:
  %
  %   in=105   puts the second control point on the far side of the chord from
  %            the first, which is what lets the curve inflect at all. A shallow
  %            `in' (150, 180) leaves both control points on the same side and
  %            the curve can only bow, however it is otherwise tuned.
  %   looseness=3  lengthens both control arms. Past about 2 they overshoot, and
  %            the overshoot is the turn back: the curve climbs again for a
  %            stretch instead of merely flattening. The same angles give a
  %            smooth S with no reversal at about 1.2, and a plain bow below 1.
  %
  % So the shape is in the tangents, not in a decoration -- a snake would repeat
  % the wobble, which reads as noise rather than one deliberate turn.
  \draw[soft, shorten >=2pt]
        ($ (slope.base west)!0.20!(slope.base east) + (0,-0.15) $)
        to[out=-85, in=105, looseness=3.0] ($(A)!0.50!(B)$);

  % --- construction levels ---------------------------------------------------
  \draw[level] (\xleft-\stub,\vz) -- (\xlev,\vz);
  \draw[level] (B) -- (\xlev,\vf);
  \draw[level] (\xtip,0) -- (\xlev,0);

  % --- dimensions ------------------------------------------------------------
  \draw[dim] (\xleft,0)  -- node[left]  {$v_0$} (\xleft,\vz);
  \draw[dim] (\xdim,0)   -- node[right] {$v_0$} (\xdim,\vz);
  \draw[dim] (\xdim,\vz) -- node[right] {$\Delta v = a\,\Delta t$} (\xdim,\vf);
  \draw[dim] (\xtot,0)   -- node[right] {$v_f = v_0 + a\,\Delta t$} (\xtot,\vf);

  % The base of the trapezoid. Set below the axis so it is read as the width
  % of the shaded region, not as a coordinate on it.
  \node[below=3pt] at (\T/2,0) {$\Delta t$};

  % --- what the area comes to ------------------------------------------------
  % An aligned, so the second line's "=" sits under the first by construction.
  % (An earlier version faked that with \phantom, on the belief that the
  % standalone wrapper had no amsmath -- it does, via unicode-math, and now
  % says so explicitly.)
  %
  % Anchored by its top-left corner rather than centred on a tuned offset. That
  % corner is where the pointer lands, so anchoring it means re-setting the
  % expression wider or taller can no longer drag the target around -- and
  % since the corner is at positive x by construction, the pointer can never be
  % left having to cross the v axis to reach it.
  \node (area) [anchor=north west] at (0.30*\T,-0.90) {$\begin{aligned}
    \Delta x &= v_0\,\Delta t + \tfrac{1}{2}(a\,\Delta t)(\Delta t) \\
             &= v_0\,\Delta t + \tfrac{1}{2}a\,(\Delta t)^2
  \end{aligned}$};

  % Leaves the area heading straight down and arrives heading down-and-left, so
  % the two tangents disagree and the curve has to bow to join them -- a single
  % sweep with no straight stretch in it. Matching the tangents (or aiming
  % either end at the other) collapses it into a diagonal line, which is the
  % one shape a pen would not make. The tail is a fraction of the figure rather
  % than a point, so it stays inside the area at any proportions.
  \draw[soft] (0.40*\T,0.70*\vz) to[out=-90, in=100, looseness=1.25]
        (area.north west);

\end{tikzpicture}
